%%%%%%%%%%%%%%%%%%%%%%%%%%%%%%%%%%%%%%%%%%%%%%%%%%%%%%%%%%%%%%%%%%%%%%%%%%%%%%%%
% VULNERABILITY FINDINGS
%%%%%%%%%%%%%%%%%%%%%%%%%%%%%%%%%%%%%%%%%%%%%%%%%%%%%%%%%%%%%%%%%%%%%%%%%%%%%%%%

\chapter[Vulnerability Findings]{\IconWarning\ Vulnerability Findings}

% Define inline badges using robust fcolorbox
\newcommand{\CriticalBadge}{\fboxsep=3pt\fboxrule=0.6pt\fcolorbox{ZWRed}{ZWRed!10}{\small\bfseries\color{ZWRed}CRITICAL}}
\newcommand{\HighBadge}{\fboxsep=3pt\fboxrule=0.6pt\fcolorbox{ZWRed}{ZWRed!10}{\small\bfseries\color{ZWRed}HIGH}}
\newcommand{\MediumBadge}{\fboxsep=3pt\fboxrule=0.6pt\fcolorbox{ZWOrange}{ZWOrange!10}{\small\bfseries\color{ZWOrange}MEDIUM}}
\newcommand{\LowBadge}{\fboxsep=3pt\fboxrule=0.6pt\fcolorbox{ZWTeal}{ZWTeal!10}{\small\bfseries\color{ZWTeal}LOW}}
\newcommand{\NegligibleBadge}{\fboxsep=3pt\fboxrule=0.6pt\fcolorbox{gray}{gray!10}{\small\bfseries\color{gray}NEGLIGIBLE}}

This chapter presents the detailed results of the vulnerability assessment performed using the \textbf{\VulScanner} Software Composition Analysis (SCA) tool. The assessment was conducted against the CycloneDX Software Bill of Materials (SBOM) generated for the \textbf{\ProjectName}.

Grype and the ZeroWatch Agent Scan compared each software component identified within the SBOM against multiple vulnerability intelligence sources, including GitHub Security Advisories (GHSA) and the National Vulnerability Database (NVD), to identify known security vulnerabilities, determine their severity, and identify available remediation guidance.

\vspace{0.4cm}

%%%%%%%%%%%%%%%%%%%%%%%%%%%%%%%%%%%%%%%%%%%%%%%%%%%%%%%%%%%%%%%%%%%%%%%%%%%%%%%%
% SCAN OVERVIEW
%%%%%%%%%%%%%%%%%%%%%%%%%%%%%%%%%%%%%%%%%%%%%%%%%%%%%%%%%%%%%%%%%%%%%%%%%%%%%%%%

\section{Scan Overview}

\begin{table}[H]
\centering
\small
\renewcommand{\arraystretch}{1.3}

\begin{tabular}{p{5cm}p{8cm}}
\toprule
\rowcolor{ZWPrimary}
\textbf{\color{white}Attribute} & \textbf{\color{white}Value} \\
\midrule
Scanner & \VulScanner \\
Assessment Type & Software Composition Analysis (SCA) \\
Target & \ProjectName \\
Input Artifact & CycloneDX SBOM \\
Scan Date & \ScanDate \\
Advisory Sources & GitHub Security Advisories (GHSA), National Vulnerability Database (NVD) \\
Total Vulnerabilities & \TotalVulnerabilities \\
Critical & \CriticalVuls \\
High & \HighVuls \\
Medium & \MediumVuls \\
Low & \LowVuls \\
Negligible & \NegligibleVuls \\
Available Fixes & Yes \\
\bottomrule
\end{tabular}
\caption{Vulnerability Scan Overview}
\end{table}

\vspace{0.4cm}

%%%%%%%%%%%%%%%%%%%%%%%%%%%%%%%%%%%%%%%%%%%%%%%%%%%%%%%%%%%%%%%%%%%%%%%%%%%%%%%%
% SCAN METHODOLOGY
%%%%%%%%%%%%%%%%%%%%%%%%%%%%%%%%%%%%%%%%%%%%%%%%%%%%%%%%%%%%%%%%%%%%%%%%%%%%%%%%

\section{Scan Methodology}

The vulnerability assessment followed the workflow below:
\begin{enumerate}
\item Generate a CycloneDX Software Bill of Materials (SBOM) using Syft.
\item Enumerate all direct and transitive software dependencies.
\item Analyze the generated SBOM using Grype and the ZeroWatch Audit Scan.
\item Compare each dependency against publicly available vulnerability databases.
\item Identify affected packages, vulnerability severity, and available fixes.
\item Generate a structured vulnerability report for security review.
\end{enumerate}

This methodology provides repeatable and automated Software Composition Analysis suitable for continuous integration and release validation.

\vspace{0.4cm}

%%%%%%%%%%%%%%%%%%%%%%%%%%%%%%%%%%%%%%%%%%%%%%%%%%%%%%%%%%%%%%%%%%%%%%%%%%%%%%%%
% SEVERITY DISTRIBUTION
%%%%%%%%%%%%%%%%%%%%%%%%%%%%%%%%%%%%%%%%%%%%%%%%%%%%%%%%%%%%%%%%%%%%%%%%%%%%%%%%

\section{Severity Distribution}

The vulnerability assessment identified a total of 1 vulnerabilities. Of these, 0 are classified as Critical, 1 as High, 0 as Medium, 0 as Low, and 0 as Negligible severity findings.

\begin{table}[H] 
\centering
\small
\renewcommand{\arraystretch}{1.3}

\begin{tabular}{p{6cm}c}
\toprule
\rowcolor{ZWPrimary}
\textbf{\color{white}Severity} & \textbf{\color{white}Count} \\
\midrule
Critical & \CriticalVuls \\
High & \HighVuls \\
Medium & \MediumVuls \\
Low & \LowVuls \\
Negligible & \NegligibleVuls \\
\bottomrule
\end{tabular}
\caption{Vulnerability Severity Distribution}
\end{table}

\vspace{0.4cm}

%%%%%%%%%%%%%%%%%%%%%%%%%%%%%%%%%%%%%%%%%%%%%%%%%%%%%%%%%%%%%%%%%%%%%%%%%%%%%%%%
% SEVERITY CHART
%%%%%%%%%%%%%%%%%%%%%%%%%%%%%%%%%%%%%%%%%%%%%%%%%%%%%%%%%%%%%%%%%%%%%%%%%%%%%%%%

\section{Vulnerability Severity Chart}

Figure~\ref{fig:vulnerability-chart} illustrates the severity distribution of all vulnerabilities identified during the assessment.

\begin{figure}[ht]
\centering

\begin{tikzpicture}
\begin{axis}[
    ybar,
    width=12cm,
    height=7cm,
    ymin=0,
    ylabel={Number of Vulnerabilities},
    symbolic x coords={Critical,High,Medium,Low,Negligible},
    xtick=data,
    xticklabel style={rotate=0,anchor=north},
    nodes near coords,
    axis x line*=bottom,
    axis y line*=left,
    bar width=25pt,
    ymajorgrids=true,
    grid style={dashed,gray!30}
]

\addplot[fill=ZWRed, draw=none] coordinates {(Critical,\CriticalVuls) (High,\HighVuls) (Medium,\MediumVuls) (Low,\LowVuls) (Negligible,\NegligibleVuls)};

\end{axis}
\end{tikzpicture}

\caption{Distribution of Vulnerabilities by Severity}
\label{fig:vulnerability-chart}
\end{figure}

%VULNERABILITY_SEVERITY_SUMMARY_TEXT%

\vspace{0.4cm}
\newpage

%%%%%%%%%%%%%%%%%%%%%%%%%%%%%%%%%%%%%%%%%%%%%%%%%%%%%%%%%%%%%%%%%%%%%%%%%%%%%%%%
% IDENTIFIED VULNERABILITIES
%%%%%%%%%%%%%%%%%%%%%%%%%%%%%%%%%%%%%%%%%%%%%%%%%%%%%%%%%%%%%%%%%%%%%%%%%%%%%%%%

\section{Identified Vulnerabilities}

\small
\renewcommand{\arraystretch}{1.3}

\begin{longtable}{p{4cm}p{2.2cm}p{2cm}p{2cm}p{3.5cm}}
\caption{Identified Vulnerabilities} \\
\toprule
\rowcolor{ZWPrimary}
\textbf{\color{white}Package} & \textbf{\color{white}Severity} & \textbf{\color{white}Installed} & \textbf{\color{white}Fixed} & \textbf{\color{white}GHSA ID} \\
\midrule
\endfirsthead

\toprule
\rowcolor{ZWPrimary}
\textbf{\color{white}Package} & \textbf{\color{white}Severity} & \textbf{\color{white}Installed} & \textbf{\color{white}Fixed} & \textbf{\color{white}GHSA ID} \\
\midrule
\endhead
actions/download-artifact & \HighBadge & v4 & 4.1.3 & GHSA-cxww-7g56-2vh6 \\
\bottomrule
\end{longtable}

\vspace{0.4cm}

%%%%%%%%%%%%%%%%%%%%%%%%%%%%%%%%%%%%%%%%%%%%%%%%%%%%%%%%%%%%%%%%%%%%%%%%%%%%%%%%
% VULNERABILITY ANALYSIS
%%%%%%%%%%%%%%%%%%%%%%%%%%%%%%%%%%%%%%%%%%%%%%%%%%%%%%%%%%%%%%%%%%%%%%%%%%%%%%%%

\section{Vulnerability Analysis}

\subsection{GHSA-cxww-7g56-2vh6}

\textbf{Affected Package} \\
actions/download-artifact

\vspace{0.2cm}
\textbf{Severity} \\
\HighBadge

\vspace{0.2cm}
\textbf{Installed Version} \\
v4

\vspace{0.2cm}
\textbf{Fixed Version} \\
4.1.3

\vspace{0.2cm}
\textbf{Description} \\
@actions/download-artifact has an Arbitrary File Write via artifact extraction

\vspace{0.2cm}
\textbf{Recommendation} \\
Upgrade to version 4.1.3 or later.

\vspace{0.5cm}
\color{ZWBorder}\hrule height 0.6pt\color{ZWDark}
\vspace{0.5cm}



%%%%%%%%%%%%%%%%%%%%%%%%%%%%%%%%%%%%%%%%%%%%%%%%%%%%%%%%%%%%%%%%%%%%%%%%%%%%%%%%
% REMEDIATION PRIORITY
%%%%%%%%%%%%%%%%%%%%%%%%%%%%%%%%%%%%%%%%%%%%%%%%%%%%%%%%%%%%%%%%%%%%%%%%%%%%%%%%

\section{Remediation Priority}

\begin{table}[H]
\centering
\small
\renewcommand{\arraystretch}{1.3}

\begin{tabular}{p{5cm}p{3cm}p{5cm}}
\toprule
\rowcolor{ZWPrimary}
\textbf{\color{white}Package} & \textbf{\color{white}Priority} & \textbf{\color{white}Recommended Action} \\
\midrule
actions/download-artifact & Immediate (Next Release) & Upgrade to version 4.1.3 \\
\bottomrule
\end{tabular}
\caption{Remediation Priority}
\end{table}

\vspace{0.4cm}

%%%%%%%%%%%%%%%%%%%%%%%%%%%%%%%%%%%%%%%%%%%%%%%%%%%%%%%%%%%%%%%%%%%%%%%%%%%%%%%%
% OVERALL RISK RATING
%%%%%%%%%%%%%%%%%%%%%%%%%%%%%%%%%%%%%%%%%%%%%%%%%%%%%%%%%%%%%%%%%%%%%%%%%%%%%%%%

\section{Overall Risk Rating}

Based on the vulnerability assessment results, the overall software supply chain risk for the \textbf{\ProjectName} is assessed as \textbf{Moderate}. Although High severity vulnerabilities were identified, they can be mitigated through package updates to available fixed versions.

\vspace{0.4cm}

%%%%%%%%%%%%%%%%%%%%%%%%%%%%%%%%%%%%%%%%%%%%%%%%%%%%%%%%%%%%%%%%%%%%%%%%%%%%%%%%
% KEY OBSERVATIONS
%%%%%%%%%%%%%%%%%%%%%%%%%%%%%%%%%%%%%%%%%%%%%%%%%%%%%%%%%%%%%%%%%%%%%%%%%%%%%%%%

\section{Key Observations}

\begin{itemize}
\item \textbf{Elevated Action Required:} 1 High vulnerabilities detected.
\item A total of 1 vulnerabilities were identified.
\item The software supply chain consists primarily of Python packages.
\item MIT is the most commonly used open-source license.
\item Continuous SBOM regeneration should accompany dependency updates.
\end{itemize}

\vspace{0.4cm}

%%%%%%%%%%%%%%%%%%%%%%%%%%%%%%%%%%%%%%%%%%%%%%%%%%%%%%%%%%%%%%%%%%%%%%%%%%%%%%%%
% SUMMARY
%%%%%%%%%%%%%%%%%%%%%%%%%%%%%%%%%%%%%%%%%%%%%%%%%%%%%%%%%%%%%%%%%%%%%%%%%%%%%%%%

\section{Summary}

The Grype and ZeroWatch Agent Audit Software Composition Analysis identified vulnerabilities affecting third-party dependencies used by the \textbf{\ProjectName}. Applying the recommended dependency updates and incorporating continuous vulnerability scanning will keep the agent secure.